\section{Elementary estimates and frequency spacing}\label{app:elementary}

We prove the sampling, derivative, and spacing estimates used below.
The spacing estimate follows the argument of de Faveri and Pandey.
Appendix~\ref{app:base} states the zeta and mean-value theorems used in
the moment bounds.

\begin{lemma}[local sampling]\label{lem:localproof}
If $p\ge1$, $f\in C^1$, and $I_x=[x-1/3,x+1/3]$, then
\[
 |f(x)|^p\ll_p\int_{I_x}(|f(u)|^p+|f'(u)|^p)\,du.
\]
Hence the same inequality can be summed over any $1$-separated set,
with a constant uniform in its cardinality.
\end{lemma}
\begin{proof}
For every $u\in I_x$, the fundamental theorem of calculus gives
$|f(x)|\le |f(u)|+\int_{I_x}|f'(v)|\,dv$. Raise this inequality to the
$p$th power, use H\"older on the integral, and average over $u$.
The intervals $I_x$ for distinct sample points have disjoint interiors.
\end{proof}

\begin{lemma}[second derivative estimates]\label{lem:secondproof}
On an interval of length at most $M$, suppose
$0<\lambda\le |g''|\le C\lambda$. Then
\[
 \left|\sum_{n\in I}e(g(n))\right|
 \ll_C M\sqrt\lambda+\lambda^{-1/2}+1.
\]
For an integral with a compactly supported $C^1$ amplitude $a$, if
$|g''|\ge\lambda>0$, then
\[
 \left|\int a(x)e(g(x))\,dx\right|
 \ll\lambda^{-1/2}(\|a\|_\infty+\|a'\|_1).
\]
\end{lemma}
\begin{proof}
The integral estimate follows by removing $|g'|\le\sqrt\lambda$, whose
length is at most $2/\sqrt\lambda$, and integrating by parts on each of
the two remaining intervals. Since $g'$ is monotone,
$\int |g''|/|g'|^2\le2/\sqrt\lambda$ there.

For the sum, the statement is trivial for $\lambda\ge1/16$. Otherwise
put $\delta=\sqrt\lambda$. The range of the monotone function $g'$
meets $O_C(M\lambda+1)$ integers. Remove the inverse images of the
$2\delta$-neighborhoods of these integers and the unit intervals meeting
their endpoints. They contain
$O_C((M\lambda+1)(\delta/\lambda+1))$ integers. On each complementary
interval the increments $g(n+1)-g(n)$ are monotone and lie between
$q+\delta$ and $q+1-\delta$ for one integer $q$. The identity
\[
 e(g(n))=\frac{e(g(n+1))-e(g(n))}{e(g(n+1)-g(n))-1}
\]
and summation by parts bound its sum by $O(\delta^{-1})$: the total
variation of the reciprocal denominator along this arc is
$O(\delta^{-1})$. There are $O_C(M\lambda+1)$ such intervals.
Adding these estimates and inserting $\delta=\sqrt\lambda$ proves the bound.
\end{proof}

\begin{lemma}[spacing, after de Faveri--Pandey]\label{lem:spacingproof}
Let $\phi(x)=\sqrt x+\eta x^{3/2}$. Uniformly for sufficiently small
$|\eta|$, $0<d<1$, and $v\in\mathbb R$,
\[
 \#\{m,n\in(N,2N]:m\ne n,\ |\phi(n/N)-\phi(m/N)-v|\le d\}
 \ll dN^2+N^{2/3}.
\]
\end{lemma}
\begin{proof}
For $d\ge c/N$, monotonicity gives $O(1+dN)$ choices for $n$ at each
$m$, which suffices. Hence suppose $d<c/N$ with $c$ small.
By interchanging $m,n$, take $v\ge0$. If any pair is counted then
$v\gg N^{-1}$ and $v\ll1$. Work on the interval where the inverse below
takes values in a fixed neighborhood of $[1,2]$, and put
$g_v(x)=\phi^{-1}(\phi(x)+v)$. Direct implicit differentiation gives
\[
 g_v'=\frac{\phi'(x)}{\phi'(g_v)},\qquad
 g_v''=\frac{\phi'(x)^2}{\phi'(g_v)}
 \left(\frac{\phi''(x)}{\phi'(x)^2}
       -\frac{\phi''(g_v)}{\phi'(g_v)^2}\right).
\]
Now $\phi'>c_1>0$, $\phi''<-c_1<0$, and
$(\phi''/\phi'^2)'$ is bounded above and below by positive constants
on this fixed interval. The last statement follows at $\eta=0$ from
$\phi''/\phi'^2=-x^{-1/2}$, and persists by continuity.
It follows that $g_v-x\asymp v$, $g_v'-1\asymp v$, and
$-g_v''\asymp v$. Each admissible $m$ must satisfy
\[
 \|f(m)\|\le C dN,\qquad f(u)=N(g_v(u/N)-u/N).
\]
Each $m$ has $O(1)$ corresponding $n$.

If $v\le N^{-1/3}$, the values $f(m)$ run through an interval of length
$O(Nv)$ and are separated by $\gg v$. Thus their number within $CdN$
of an integer is $O((1+Nv)(1+dN/v))=O(Nv+dN^2)$, as $Nv\gg1$.

For $v>N^{-1/3}$ let $b=C'dN+(v/N)^{1/3}$.
If $b$ exceeds a fixed small constant, the trivial bound suffices.
Otherwise periodize a nonnegative smooth bump at scale $b$, equal to at
least one on $[-CdN,CdN]$. Its Fourier coefficients satisfy
$|c_h|\ll_A b(1+b|h|)^{-A}$. Since $|f''|\asymp v/N$,
Lemma~\ref{lem:secondproof} and Fourier expansion give
\[
 \#\{m:\|f(m)\|\le CdN\}
 \ll Nb+\sqrt{Nv}\,b^{-1/2}+\sqrt{N/v}\,b^{1/2}+1.
\]
Here $Nb\ll dN^2+N^{2/3}$,
$\sqrt{Nv/b}\ll N^{2/3}v^{1/3}\ll N^{2/3}$, and
$\sqrt{Nb/v}\ll\sqrt{N/v}\ll N^{2/3}$.
The absolutely convergent smooth Fourier series and the proved
second-derivative bound establish this spacing estimate.
\end{proof}

\section{The curved multiplier and all even energies}\label{app:energyproof}

\begin{proposition}\label{prop:multproof}
Lemma~\ref{lem:mult} holds with an ordinary implied constant in place of
$\prec$. The constants are uniform on fixed support enlargements where
$|\kappa|W$ is sufficiently small.
\end{proposition}
\begin{proof}
Write $\lambda=\xi_m-\xi_n$, $\nu=\chi_m-\chi_n$,
$R_\nu=1+B|\nu|$ and $a_\nu(u)=\beta(u/W)e(\nu r_0(u))$.
On this support,
\[
 |W^j r_0^{(j)}(u)|\ll_j B,\qquad
 r_0''(u)=-2\eta\kappa(1+\kappa u)^{-3},\qquad
 |\nu|\asymp|\lambda|.
\]
The integral part of Lemma~\ref{lem:secondproof}, and the trivial estimate
when $B|\nu|\le1$, show $|\widehat a_\nu(v)|\ll WR_\nu^{-1/2}$.
If $W|v|>C R_\nu$, rescale $u=Wx$ and integrate with
$(2\pi i(\nu W r_0'(Wx)-Wv))^{-1}d/dx$. Its denominator is comparable
to $W|v|$, and every normalized derivative of the nonlinear part is
$O_j(B|\nu|)$. Repeating gives
$|\widehat a_\nu(v)|\ll_J W(1+W|v|)^{-J}$ in this region.

Lemma~\ref{lem:spacingproof} puts $O(N^2/W)$ differences $\lambda$ in
each interval of length $1/W$, since $W\le N^{4/3}$.
On $|u-\lambda|\le CR_\nu/W$, the smallness of $B/W$ gives
\[
 1+B|\lambda|\asymp1+B|u|.
\]
For example $B||\lambda|-|u||\le C(B/W)(1+B|\lambda|)$ and this term
can be absorbed in both directions. The near region has
$O(1+B|u|)$ cells. Multiply their population by
$W(1+B|u|)^{-1/2}$. The result is $O(N^2\sqrt{1+B|u|})$.
On the complement use the separate far estimate just proved, summing
$O(N^2/W)W(1+|j|)^{-J}$ over cells at distance $j/W$ from $u$.
This contributes $O(N^2)$ and proves the statement.
\end{proof}

\begin{proposition}[even energy extension]\label{prop:allenergy}
Under the hypotheses of the curved model, for every fixed integer $j\ge2$
and $1\le D\le W$,
\[
 E_{2j}(S;D)\ll_j
 \left(\frac NZ\right)^{4j-4}
 \left(1+\frac BD\right)^{j-1}E_2(S;D)
 +\frac{ND}{Z^2}r^{2j-1}.
\]
The implicit constants depend on the fixed cutoffs and on $j$.
\end{proposition}
In particular this proves Proposition~\ref{prop:curv}.
\begin{proof}
Set $\mathcal F(u)=\sum_{w\in S}e(wu)$. For each ordered $(2j-1)$-tuple put
$\tau=w_1+\cdots+w_j-w_{j+1}-\cdots-w_{2j-1}$.
Multiply its last-point count by $Z^2$ and use
Lemma~\ref{lem:localproof} on each eligible last point.
One fixed nonnegative $g\in C_c^\infty$, equal to one on $[-2,2]$,
majorizes the resulting local intervals by $g((\tau-u)/D)$.
Since
\[
 D_h'(u)=2\pi i\left[\sum a_n\xi_ne(u\xi_n+r_0(u)\chi_n)
 +r_0'(u)\sum a_n\chi_ne(u\xi_n+r_0(u)\chi_n)\right]
\]
and $|r_0'|\ll B/W\ll1$, we obtain three integrals with fixed
coefficient sequences bounded by an absolute constant.
Before enlargement all integration points lie in $[-1,W+1]$, because
they are local neighborhoods of points of $S$. Insert a fixed
nonnegative cutoff $\beta(u/W)$ equal to one there, supported in
$[-3W,3W]$, and then enlarge the integrals to the whole line.
Each coefficient class and the cutoff remain fixed across tuples and $j$.

For each coefficient class the diagonal after opening the square is
$O_j(NDr^{2j-1})$. For an off-diagonal frequency pair, convolution and
Fourier inversion give the exact identity
\[
 \int g((\tau-u)/D)a_\nu(u)e(\lambda u)\,du
 =D\int\widehat g(Dv)\widehat a_\nu(v-\lambda)e(\tau v)\,dv.
\]
Summing over the $2j-1$ points gives $\mathcal F(v)|\mathcal F(v)|^{2j-2}$.
Proposition~\ref{prop:multproof} therefore bounds the absolute
off-diagonal contribution by
\[
 CN^2\sqrt{1+B/D}\,D\int\Psi(Dv)|\mathcal F(v)|^{2j-1}\,dv,
 \qquad \Psi(v)=(1+v^2)^{-J},
\]
with $J$ fixed sufficiently large. We used
$\sqrt{1+B|v|}\le\sqrt{1+B/D}\sqrt{1+D|v|}$ and rapid decay of
$\widehat g$.

For each fixed $h\ge1$, expansion and Fourier inversion show
\[
 D\int\Psi(Dv)|\mathcal F(v)|^{2h}\,dv\ll_h E_{2h}(S;D).
\]
For the tail bound, $|\widehat\Psi(u)|\ll_A(1+|u|)^{-A}$ follows by integration
by parts. Split the tuple discrepancies into dyadic shells and use
Lemma~\ref{lem:enlarge} on the multiset of $h$-fold sums. The shell at
radius $2^aD$ contributes $O(2^{-(A-1)a}E_{2h}(S;D))$.
H\"older with exponents $2j-2$ and $(2j-2)/(2j-3)$ now yields
\[
 Z^2E_{2j}\ll_j NDr^{2j-1}
 +N^2\sqrt{1+B/D}\,
 E_2^{1/(2j-2)}E_{2j}^{(2j-3)/(2j-2)}.
\]
If the first term is at least half the left side, it gives the stated
diagonal bound. Otherwise divide by the last power of $E_{2j}$ and
raise to $2j-2$. This proves the statement, including all ordered tuples.
\end{proof}

\section{Exact oscillatory integrals and their parameters}\label{app:phaseproof}

\begin{lemma}[two pairs with finitely many derivatives]\label{lem:smoothpairs}
For a fixed compact interval and $f\in C^2$ on a slightly larger interval,
\[
 |f(v_1)+f(v_2)-f(v_3)-f(v_4)|
 \ll\|f\|_{C^2}(|P_v|+|Q_v|).
\]
The statement applies to any fixed parameter derivative of $f$ provided
that derivative has a bounded $C^2$ norm.
\end{lemma}
\begin{proof}
Write a pair as $s/2\pm r$ and put $\rho=r^2$.
For $G(s,\rho)=f(s/2+r)+f(s/2-r)$,
\[
 G_s=\tfrac12(f'(s/2+r)+f'(s/2-r)),\quad
 G_\rho=\tfrac12\int_{-1}^1 f''(s/2+vr)\,dv.
\]
These formulas include $r=0$ by continuity. Extend $f$ smoothly to a
larger fixed interval if necessary and use the segment between the two
points $(s,\rho)$ in the half-plane $\rho\ge0$.
For the two pairs, $s-s'=P_v$ and
$\rho-\rho'=Q_v/2-(s^2-s'^2)/4$.
The mean value theorem proves the estimate.
Applying the same argument to a parameter derivative proves uniformity.
\end{proof}

\begin{lemma}[exact Morse and Fresnel identity]\label{lem:exactmorse}
Let $p$ range in a fixed compact parameter set. Suppose $f(p,x)$ has,
on a fixed interval containing the support of a smooth amplitude $a$,
a unique critical point $x_0(p)$, and $s f_{xx}\ge c>0$ there for one
$s\in\{-1,1\}$. Suppose also that $f$ and $a$ have bounded derivatives
of all fixed orders under consideration. If $Y\ge1$, then exactly
\[
 \int a(p,x)e(Yf(p,x))\,dx
 =Y^{-1/2}e(Yf(p,x_0(p)))\,b(p,Y),
\]
where $(Y\partial_Y)^h\partial_p^\alpha b\ll_{h,\alpha}1$.
The support of the amplitude is allowed to contain $x_0$ or lie away from it.
\end{lemma}
\begin{proof}
The implicit function theorem and differentiation of $f_x(p,x_0)=0$
give bounded derivatives of $x_0$. Taylor's integral formula gives
\[
 v=(x-x_0)\left(s\int_0^1(1-u)f_{xx}(p,x_0+u(x-x_0))\,du\right)^{1/2}.
\]
It has bounded derivatives and an inverse with bounded derivatives;
the identity $2v v_x=s f_x$ and the lower curvature bound show
$v_x\asymp1$, also at $x_0$. The phase becomes
$f(p,x_0)+s v^2$. The transformed amplitude $c(p,v)$ is smooth and
compactly supported in a fixed interval, also after extension by zero.
Fourier inversion and the Gaussian integral give
\[
 b(p,Y)=\frac{e(s/8)}{\sqrt2}
 \int_{\mathbb R}\widehat c(p,\xi)e(-s\xi^2/(4Y))\,d\xi.
\]
One may first insert $e^{-\epsilon v^2}$, evaluate the ordinary Gaussian
integral, and pass to the limit by Schwartz domination on the Fourier
side; this justifies the oscillatory Gaussian identity.
Each $Y\partial_Y$ inserts a polynomial in $\xi^2/Y$.
Parameter derivatives fall on $\widehat c$. Its Schwartz bounds prove
all the stated estimates. The construction uses an interval containing both the support and $x_0$.
It also covers the case $a(x_0)=0$.
\end{proof}

\begin{proposition}[critical value and tuple derivatives]\label{prop:exactphaseproof}
Lemma~\ref{lem:phase} and \eqref{eq:psi} hold. More explicitly, if
$L=NM$, $t=m/M$ is in a fixed interval away from zero, $v_i=w_i/L$
are bounded and $\gamma=\kappa L$, then
\[
 \Phi(t,x)=\frac Q{4Lt}-\frac{Px}{2}+C(t,x),\qquad
 |\partial_t^a\partial_x^b C|\ll_{a,b}|\eta|(|P|+F),
\]
and
\[
 C_{xx}=\frac{3\eta Q}{2Lt}-3\eta Px+E(t,x),\qquad
 |\partial_t^a\partial_x^b E|
 \ll_{a,b}(|\gamma|+|\eta|)|\eta|(|P|+F).
\]
All estimates are on fixed compact $x$ intervals, with their smallness
constants chosen after those intervals.
\end{proposition}
\begin{proof}
The pair phase before Poisson integration is
\[
 y\{\Delta u^{1/2}+\eta_{\rm eff}\Delta u^{3/2}-u\},\quad
 y=Nm,\quad \Delta=(z-w)/y.
\]
Writing $r=\sqrt u$ in its critical equation gives
$3\eta_{\rm eff}\Delta r^2-2r+\Delta=0$.
The root remaining in a fixed positive interval as $\eta\to0$ is
$r=\Delta/(1+\sqrt{1-3\eta_{\rm eff}\Delta^2})$.
The other root is outside the fixed integration support. On the
stationary range $\Delta$ is positive and bounded above and below,
and $|\eta_{\rm eff}\Delta^2|<1/4$ follows from smallness of $|\eta|$.
Substitution, using $\frac34\tau\Omega^2-\Omega+1=0$, gives exactly
\[
 \Theta=\frac{y\Delta^2}{12}\Omega(4-\Omega).
\]
The analytic expansion at $\tau=0$ is
\[
 \Theta=\frac{y\Delta^2}{4}
       +\frac{y\eta_{\rm eff}\Delta^4}{8}
       +y\eta_{\rm eff}^2\Delta^6\Psi(\tau),
 \qquad \Psi(0)=\frac9{64}.
\]
This identity keeps the full analytic remainder function in the phase.

At $z=Ltx$ its leading term sums to $Q/(4Lt)-Px/2$.
The next term sums to
\[
 \frac{\eta L}{8t^3(1+\gamma tx)}
 \sum_i\epsilon_i\frac{(tx-v_i)^4}{1+\gamma v_i}.
\]
The remainder is $\eta^2L$ times the signed two-pair difference of a
uniformly smooth function of $v_i$, also after any fixed $t,x$
derivatives. Lemma~\ref{lem:smoothpairs} bounds it by
$O_{a,b}(\eta^2(|P|+F))$. The same lemma bounds the quartic term by
$O_{a,b}(|\eta|(|P|+F))$.
Replacing $\gamma$ by zero in its second $x$ derivative costs
$O_{a,b}(|\gamma\eta|(|P|+F))$ by the mean value theorem in $\gamma$
and the two-pair estimate. At $\gamma=0$ direct differentiation gives
$3\eta Q/(2Lt)-3\eta Px$. This proves every claimed derivative bound.
\end{proof}

\begin{proposition}[tuple Morse formula]\label{prop:tuplemorseproof}
Lemma~\ref{lem:morse} holds. The nonstationary and small-$H$ statements
used in Proposition~\ref{prop:kernel} also hold with bounded derivatives
of every fixed required order.
\end{proposition}
\begin{proof}
Let $H=|\eta|F$. When $|P|>C(1+H)$, the first $x$ derivative has
magnitude $\asymp|P|$. Factor $e(Q/(4Lt))$ and integrate by parts in
$x$. Every fixed $t$ derivative costs at most a fixed power of
$|P|+H$; choosing that many extra integrations gives
$O_{a,J}((1+|P|)^{-J})$ for the remaining amplitude and its derivatives.
For $H\le1$ and $|P|\le C(1+H)$ all derivatives of $\psi$ are bounded,
so $e(\psi)$ is a bounded smooth amplitude.

Suppose $H\ge1$ and $|P|\le CH$. The preceding proposition gives
\[
 \psi_{xx}/H=\operatorname{sgn}(\eta Q)\frac3{2t}
             +O(|\eta|+|\kappa|L)
\]
for positive $t$; the sign is adjusted for negative $t$.
Moreover $\psi_x/H$ is an affine function of $x$ with bounded intercept
and slope separated from zero, plus $O(|\eta|+|\kappa|L)$ in every
fixed derivative on a fixed interval. For the intercept one uses
$|\sum\epsilon_iw_i^3|\ll L|Q|+L^2|P|$, from
Lemma~\ref{lem:smoothpairs}. Choose a fixed larger interval containing
the amplitude support and all roots of these affine functions, then
choose the small constants. The derivative has opposite signs at the
ends, so there is exactly one critical point and its derivatives are
bounded. Apply Lemma~\ref{lem:exactmorse} with $Y=H$ and $f=\psi/H$.
Here $H$ is fixed as $t$ varies. It yields
$H^{-1/2}b(t)e(Q/(4Lt)+\psi(t,x_0(t)))$, with
$b^{(a)}\ll_a1$ and $(d/dt)^a\psi(t,x_0(t))\ll_aH$.
The construction includes critical points outside the amplitude support.
\end{proof}

\section{A complete signed duality calculation}\label{app:dualityproof}

Put $L=NM$. We write the bound as a sum over height blocks. Partition $S$ into
intervals of length $CL$, independently of $m$, where $C$ is fixed
sufficiently large. For each outer block $S_j$ and neighboring inner
block $S_{j+h}$, $h\in\{-1,0,1\}$, translate by its fixed origin $u_j$.
Let $\mathcal B_{M,j,h}$ be the translated inner set. Then
$|w|\ll L$ in $\mathcal B_{M,j,h}$. The parameters become
\[
 \eta_j=\eta(1+\kappa u_j)^{-2},\qquad
 \kappa_j=\kappa(1+\kappa u_j)^{-1}.
\]
The common phase at $u_j$ cancels in the pair difference, since
$h(u_j+v)-h(u_j)=\eta_jv/(1+\kappa_jv)$. Thus the translation can be
performed after the global duality step with its coefficients preserved.

\begin{proposition}[signed blocks]\label{prop:dualityproof}
There are smooth amplitudes in the two fixed Sobolev classes and in
each sign of $m$ for which
\begin{equation}\label{eq:preciseduality}
 r^2Z^2\ll N^2r+
 \sum_{M,h,\pm,\ell}N^{7/4}M^{1/2}
 \sum_j |S_j|^{3/4}\mathcal T_{M,j,h,\pm,\ell}^{1/4},
\end{equation}
where $\ell=0,1$, $\mathcal T\ge0$, and
\[
 \mathcal T_{M,j,h,\pm,\ell}
 =\sum_{\mathbf w\in\mathcal B_{M,j,h}^4}
 b_{w_1}b_{w_2}\overline{b_{w_3}b_{w_4}}
 \sum_{m\text{ in the fixed sign}}\int A_{\ell,\mathbf w}(m/M,x)
 e(\Phi(m/M,x))\,dx.
\]
All normalized derivatives of the amplitudes are bounded. In particular
$\mathcal T\le\sum_{\mathbf w}|K_{M,j,h,\pm,\ell}(\mathbf w)|$.
The number of neighbors and signs is bounded independently of $N,W,r$.
\end{proposition}
\begin{proof}
Choose $b_w$ of modulus one so that $b_wD_h(w)=|D_h(w)|$.
Cauchy's inequality only in the $n$ variable gives
\[
 r^2Z^2\le N\sum_n\varphi(n/N)
 \left|\sum_{w\in S}b_w
 e(w\sqrt{n/N}+h(w)(n/N)^{3/2})\right|^2,
\]
where $\varphi\ge1$ on $[1,2]$ and has fixed compact positive support.
Poisson summation changes the right side exactly to
\[
 N^2\sum_{z,w\in S}b_z\overline b_w\sum_{m\in\mathbb Z}
 \int\varphi(u)e((z-w)\sqrt u+(h(z)-h(w))u^{3/2}-Nm u)\,du.
\]
Poisson follows, for this smooth compact support, by periodizing the
summand and taking its absolutely convergent Fourier series.
For $m=0$ the derivative is comparable to $z-w$ unless $z=w$;
integration by parts and $1$-spacing give $O(N^2r)$.

For $m\ne0$, stationary support requires $\Delta=(z-w)/(Nm)$ in a
fixed compact subinterval of $(0,\infty)$. Choose a smooth cutoff in
$\Delta$ equal to one on a neighborhood of all such ratios. Off that
cutoff, repeated integration by parts bounds the integral by
$O_A((1+|z-w|+N|m|)^{-A})$.
For each fixed $z$, $1$-spacing implies
$\sum_w(1+|z-w|+N|m|)^{-A}\ll(1+N|m|)^{1-A}$.
Summing in $m$ bounds this part by $O(N^2r)$.

On the cutoff use the exact identity
$h(z)-h(w)=\eta(z-w)/((1+\kappa z)(1+\kappa w))$.
The pair phase divided by $|Nm|$ has one nondegenerate critical point
in a fixed positive enlargement of the $u$ support. Its second
derivative has constant sign there, bounded away from zero. Apply
Lemma~\ref{lem:exactmorse} with $Y=|Nm|$.
The integral is \emph{exactly}
\[
 (N|m|)^{-1/2}a_m(z,w)e(\Theta_m(z,w)).
\]
This exact representation keeps the full stationary integral. Smooth cutoffs in
$\Delta$ are included in $a_m$. Its normalized derivatives in
$m/M,z/L,w/L$ are bounded: derivatives of $Y$ are controlled by the
$Y\partial_Y$ estimates, and derivatives of the remaining parameters
cost bounded powers of $|\kappa|L$.

Fix one dyadic scale, neighbor and sign. The contribution is bounded by
\[
 N^{3/2}M^{-1/2}
 \sum_j\sum_m\omega(m/M)\sum_{z\in S_j}|f_{m,0}(z-u_j)|,
\]
where $\omega$ is a fixed nonnegative smooth dyadic weight and
\[
 f_{m,0}(z)=\sum_{w\in\mathcal B_{M,j,h}}b_w a_{m,0}(z,w)e(\Theta_m(z,w)).
\]
Weighted H\"older in $(m,z)$ gives the factor
$(M|S_j|)^{3/4}$ and the fourth root of
\[
\sum_m\omega(m/M)\sum_{z\in S_j-u_j}|f_{m,0}(z)|^4.
\]
By Lemma~\ref{lem:localproof} the latter is bounded by
\[
 \sum_m\omega(m/M)\int\chi_j(z/L)
 (|f_{m,0}(z)|^4+|f_{m,1}(z)|^4)\,dz,
 \quad f_{m,1}=\partial_zf_{m,0}.
\]
The derivative amplitude is
$a_{m,1}=\partial_z a_{m,0}+2\pi i(\partial_z\Theta_m)a_{m,0}$.
At the critical point the envelope identity gives
$\partial_z\Theta_m=\sqrt{u_0}+h'(z)u_0^{3/2}=O(1)$,
with all normalized derivatives bounded. Thus both classes qualify.

Set $z=Nmx$. The Jacobian is $N|m|=L|m|/M$.
Its smooth factor $|m|/M$, and the cutoff $\chi_j((m/M)x)$, belong in
$A_{\ell,\mathbf w}$. Opening each fourth power before taking absolute
values produces exactly $\mathcal T$ above. It is nonnegative because
it was the integral of a fourth power against nonnegative weights.
The remaining external factor is
\[
 N^{3/2}M^{-1/2}(M|S_j|)^{3/4}L^{1/4}
 =N^{7/4}M^{1/2}|S_j|^{3/4}.
\]
Both $b_w$ and the translation remain fixed across $m$. This proves
\eqref{eq:preciseduality} and accounts for both signs, both derivative
classes, all neighboring blocks and all nonstationary modes.
\end{proof}

\begin{corollary}[the energy bound]\label{cor:blockconsumer}
Define
\[
 \mathfrak K_M=\sum_j\sum_{\mathbf w\in\mathcal B_{M,j,h}^4}
 |K_{M,j,h,\pm,\ell}(\mathbf w)|,
\]
separately for each finite class, and
$\cJ=\max_{M,h,\pm,\ell}(M/\Mst)^3M^{-1}\mathfrak K_M$.
Then \eqref{eq:rJ} and \eqref{eq:Jbound} hold with this definition.
\end{corollary}
\begin{proof}
H\"older in $j$ bounds the block sum in
\eqref{eq:preciseduality} by $r^{3/4}\mathfrak K_M^{1/4}$.
Dyadic summation loses only $O(\log N)$.
For each fixed $h$, the inner blocks partition $S$; hence
\[
 \sum_jE_4(S_{j+h};D)\le E_4(S;D),\qquad
 \sum_j|S_{j+h}|E_2(S_{j+h};H_0)\le rE_2(S;H_0).
\]
The discrepancies $P$ are translation invariant. The local $Q$, hence
$F$, is calculated after its block's fixed translation; the small-$F$
proof is applied there before going back to the translation invariant
close-pair counts. Each block keeps its own translated $Q$.
The pointwise kernel bounds are uniform over the translated parameters,
so energy bounds can be applied after summing the block tuple counts.
The remainder of the calculation is exactly the algebra in
Proposition~\ref{prop:model}, with $\mathfrak K_M$ in place of a single
block kernel. This proves both stated inequalities.
\end{proof}

\section{The reduction to the zeta function}\label{app:transferproof}

\begin{lemma}[partial summation and interval completion]\label{lem:completionproof}
The first phase truncation in Lemma~\ref{lem:atkinson} changes a partial
sum by at most a constant times the maximum of the partial sums with
the truncated phase. For finitely many blocks and any phase $p_j(w,n)$,
the sum over sample points of these maxima is bounded by
$C\log(2N)$ times the sum of full sums after multiplying every
coefficient by one common factor $e(hn/(2N))$ for a suitable integer $h$.
\end{lemma}
\begin{proof}
For $A(x)=\sum_{N<n\le x}c_n$ and $v(n)=e(r_t(n))$, Abel summation gives
\[
 \sum_{N<n\le y}c_nv(n)=A(y)v(y)-\int_N^y A(x)v'(x)\,dx.
\]
The factor is at most $1+2\pi\int_N^{2N}|r_t'(x)|dx=O(1)$ when
$N^{5/2}T^{-3/2}\ll1$.
For interval completion, expand the indicator of $1\le n-N\le\Delta$
on $\mathbb Z/(2N)\mathbb Z$. Its coefficient at $h$, $-N<h\le N$,
is bounded by $C/(1+|h|)$, directly from the finite geometric sum.
Thus each maximum is bounded by
\[
 C\sum_{-N<h\le N}\frac1{1+|h|}
 \left|\sum_{n\sim N}a_{j,n}e(hn/(2N))e(p_j(w,n))\right|.
\]
Sum this inequality over the blocks and sample points. Choosing the
$h$ with the largest resulting sum proves the statement.
This argument applies to every phase $p_j(w,n)$. Nonintegral endpoints are
handled by their integer parts or by a cyclic group of length $\asymp N$.
\end{proof}

\begin{proposition}[refined pair count from sparse fourth moments]
\label{prop:pairproof}
Assume the sparse fourth-moment input \eqref{eq:basefourth} and the
global fourth moment. Let $t_1,\ldots,t_R$ be $2G$-separated, and assume
\[
\int_{t_j-G}^{t_j+G}|\zeta(1/2+iu)|^4du\ge A\ge1.
\]
Then for all $H\ge1$,
\[
 E_2(\{t_j\};H)\prec R+RH/A+T\sqrt R/A^{3/2}.
\]
\end{proposition}
\begin{proof}
Disjointness and the fourth moment give $RA\prec T$.
For $H<G$ there are only diagonal pairs; for $H\ge T$ use
$R^2\prec RT/A$. Otherwise partition into cells of length $H$ with
occupancies $m_j$. Adjacent-cell Cauchy gives $E_2\le3\sum_jm_j^2$.
It suffices to treat cells of occupancy in $[m,2m)$, their number being
$J$, since there are $O(\log T)$ scales. We have $Jm\le R$.
Fix a small $\delta>0$.
If $mA\le T^\delta\max(H,\sqrt T)$, then
\[
 Jm^2\le Rm\ll T^\delta(RH/A+R\sqrt T/A)
 \prec T^\delta(RH/A+T\sqrt R/A^{3/2}).
\]
In the other case set $L=\min(mAT^{-\delta},T)>\max(H,\sqrt T)$.
The $G$-neighborhoods of the points in these $J$ cells can be covered
by $O(J)$ intervals of length $L$, with a bounded number of disjoint
families. Endpoint pieces can be handled on enlarged dyadic height
intervals. Therefore the sparse fourth moment gives
\[
 JmA\prec JL+T\sqrt{J/L}.
\]
Choose its arbitrary-small-power loss smaller than $\delta/2$ to
absorb $JL$. Squaring and using $mA\prec T$ gives
$Jm^3\prec T^{2+O(\delta)}A^{-3}$.
Thus $Jm^2=\sqrt{Jm}\sqrt{Jm^3}\prec
T^{1+O(\delta)}\sqrt R A^{-3/2}$.
As $\delta$ is arbitrary the proposition follows.
\end{proof}

\begin{proposition}[clustering from the stated zeta theorems]
\label{prop:clusterproof}
Proposition~\ref{in:dfp}, with $1/10<\sigma\le1/8$, follows from
\eqref{eq:baselocal}, \eqref{eq:baseatkinson},
\eqref{eq:basefourth} and the global fourth moment. The same conclusion
holds for $1/10<\sigma<1/6$ when the corresponding parameter conditions
for the curved model are satisfied.
\end{proposition}
\begin{proof}
Start with any $1$-separated zeta large-value set of size $R$, and put
$G=V^2(\log T)^{-100}$. Partition into cells of length $2G/3$.
Dyadic occupancy selection and one residue class modulo four give
a $2G$-separated set $\mathcal S$ of representatives, of size
$\succ R/F_c$, each representing $\asymp F_c\ge1$ original points.
By the local inequality \eqref{eq:baselocal}, their $G$-neighborhoods
carry second and fourth moments at least $\succ F_cV^2$ and
$\succ F_cV^4$, respectively.

The term $G\log T$ in \eqref{eq:baseatkinson} is smaller than the
second-moment lower bound. Select one dyadic $N$ for at least a
logarithmic fraction of the representatives. Then
\[
 T^{1/3}\ll N\prec T/V^4,\qquad
 \max_\Delta|S_t(\Delta,N)|\succ F_c(TN)^{1/4}.
\]
Since $\sigma>1/10$, the strict margin gives $N<T^{3/5}$ after all
logarithmic factors. Divide the ordinates into blocks of length
$\asymp W\sqrt{T/N}$. There are $O(1+\sqrt{TN}/W)$ blocks.
The coordinate in Lemma~\ref{lem:atkinson} sends each to $[0,CW]$;
fixed constants in $W$ may be chosen to obtain $[0,W]$.
The exact rational height increment, partial summation, the divisor
bound and Lemma~\ref{lem:completionproof} give a sum of absolute full
model values at least $\succ R(TN)^{1/4}$.

Discard values below a sufficiently small arbitrary-small-power
multiple of $F_c(TN)^{1/4}$; their total is less than half this last
sum. Dyadic selection among the remaining values gives $Z$ and sets
of total size $R_1$ with $Z\succ F_c(TN)^{1/4}$ and
$ZR_1\succ R(TN)^{1/4}$. This argument uses the summed inequality after selecting one common
Fourier index; the set and its level $Z$ are then selected together.

The transformed spacing is $\gg G\sqrt{N/T}$. Color the points of
each block in their order into $O(1+G^{-1}\sqrt{T/N})$ classes so each
class is $1$-separated. Because $N\prec T/V^4$ and $W\le\sqrt{TN}$,
the total number of classes is $\prec T/(V^2W)$.
Each point occurs once. For two points in a single class,
$|w-w'|\le D$ implies $|t-t'|\ll D\sqrt{T/N}$.
Apply Proposition~\ref{prop:pairproof} to the entire selected set of
ordinates $t$ with $A\succ F_cV^4$.
This proves \eqref{eq:paircount} for every $D\ge1$, including its
large-radius regime. The $1$-spacing bound is independent and remains
available simultaneously.
\end{proof}

\section{The upper-amplitude range}
\label{app:upperproof}

\begin{proposition}\label{prop:curvatureline}
Subject to the zeta inputs of Proposition~\ref{prop:clusterproof},
\[
 \cN(T,T^\sigma)\prec T^{(94-500\sigma)/67},
 \qquad 1/8\le\sigma\le5/38.
\]
\end{proposition}
This statement follows from the direct oscillatory-integral estimates.
\begin{proof}
Proposition~\ref{prop:tuplemorseproof}, applied at each fixed $m$, gives
$|I_m|\ll D^{-1/2}(1+D/U)^{-A}$ for every $M$ and every tuple.
Use this pointwise bound on all tuples, including the small-$F$ range. The signed-block calculation and the curved energy bound give
\[
 r\prec N^2/Z^2+N^{8/5}W^{3/5}Z^{-12/5}A_0(S)^{1/5}
 +N^{5/2}W^{3/2}Z^{-5}U^{1/4}
 +B^{1/4}N^2W^{3/4}Z^{-3}.
\]
The weighted pair argument with weight $D^{-1/2}$ is valid also for
$k=0$. Thus the reduction to the zeta function gives exactly the terms
$B_1,\ldots,B_6$ with
$\lambda=\mu=1$ and $k=0$. Put
\[
 \omega(\sigma)=\frac{36+48\sigma}{67},\qquad
 q=\frac{\omega(\sigma)}{1-4\sigma},\qquad W=N^q.
\]
Here $84/67\le q\le4/3$. All six exponents increase with
$\nu=\log_TN$. At $\nu=1-4\sigma$, terms 4 and 5 equal
$E_c=(94-500\sigma)/67$. The margins for terms 1, 2, 3 and 6 are
\[
 \frac{4(21\sigma-1)}{67},\quad
 \frac{5-38\sigma}{134},\quad
 \frac{9-55\sigma}{67},\quad
 \frac{5-38\sigma}{268},
\]
which are nonnegative throughout the closed interval.
For the model ranges and pair absorption the exponents of $T$ in
$|\kappa|W$, $|\eta|N^2/W$, $|\eta|N^{1/8}$,
$N^{5/2}T^{-3/2}$ and $W\sqrt{N/T}V^{-4}$ are, respectively,
\[
 \frac{-31+182\sigma}{67},\quad
 \frac{98-852\sigma}{67},\quad
 \frac18-\frac92\sigma,\quad
 1-10\sigma,\quad
 \frac{36-354\sigma}{67}.
\]
They are strictly negative at both endpoints. These quantities also
have positive $\nu$ slopes for fixed $q$, except the separately treated
$\sqrt{T/N}$ term, which is at most $T^{1/3}<V^4$.
The endpoint $q=4/3$ is permitted in the spacing estimate.
This proves the proposition. Its contribution to $M_6$ is at most
$T^{327/268+\varepsilon}$, attained at $\sigma=1/8$.
\end{proof}

\section{Local zeta values and consequences of sparse fourth moments}\label{app:localzeta}

\begin{proposition}[local values from harmonic measure]\label{prop:localzetaproof}
For every fixed $p>0$,
\begin{equation}\label{eq:baselocal}
 |\zeta(1/2+it)|^p\ll_p
 1+\log(2+|t|)\int_{t-1/3}^{t+1/3}|\zeta(1/2+iu)|^p\,du.
\end{equation}
\end{proposition}
\begin{proof}
We give the localization argument, including the analytic estimates it
requires. The elementary functional equation gives, uniformly for
$0\le u\le1/\log T$ and $\tau\asymp T$,
\begin{equation}\label{eq:localreflection}
 |\zeta(1/2-u+i\tau)|\ll |\zeta(1/2+u+i\tau)|.
\end{equation}
Here is one derivation of the facts needed in this statement. Periodizing
the Gaussian, whose Fourier transform is $x^{-1/2}e^{-\pi\xi^2/x}$,
gives $\theta(x)=x^{-1/2}\theta(1/x)$. Mellin integration, initially
for $\Re s>1$, and splitting at $x=1$, then gives
\[
 \pi^{-s/2}\Gamma(s/2)\zeta(s)
 =-\frac1s-\frac1{1-s}
 +\frac12\int_1^\infty(\theta(x)-1)
       \bigl(x^{s/2}+x^{(1-s)/2}\bigr)\,\frac{dx}{x}.
\]
The right side continues meromorphically and is unchanged by
$s\mapsto1-s$. The gamma quotient therefore has modulus
$\asymp(1+|\tau|)^u$ at $s=1/2-u+i\tau$, which proves
\eqref{eq:localreflection}. The gamma estimate follows
by applying Euler summation twice to the logarithmic gamma product:
on $\Re z>0$ it gives
\[
 \log\Gamma(z)=(z-\tfrac12)\log z-z+c+\frac1{12z}
 -\frac12\int_0^\infty\frac{B_2(\{x\})}{(x+z)^2}\,dx.
\]
The integral is $O(|\Im z|^{-1})$ when $\Re z$ is in a fixed positive
compact interval; take real parts in the two gamma factors.
Conjugation accounts for the sign of $\tau$. Also Euler summation in
the defining Dirichlet series, with a cutoff $X\asymp1+|\tau|$, gives
\[
 \zeta(s)=\sum_{n\le X}n^{-s}+\frac{X^{1-s}}{s-1}
          +O(X^{-\Re s}+|s|X^{-\Re s}/\Re s)
\]
for $\Re s>0$ away from $s=1$. It gives the polynomial growth used
below directly from Euler summation.

We prove the bound for $T\le t\le2T$ with $T\ge10$. On bounded intervals, continuity gives the result.
Set $s_0=1/2+it$, $\delta=1/\log T$, and choose a fixed sufficiently
large $C=C_p$. On the strip $1/2\le\Re s\le3/2$ put
\[
 G(s)=\frac{s-1}{s+1}\zeta(s)
 \exp\bigl(C\log T\bigl[(s-s_0)^2-(s-s_0)\bigr]\bigr).
\]
The pole is canceled. On either boundary the exponential has modulus
$\exp(-C\log T\,|\Im s-t|^2)$. On the right boundary $G$ is bounded
by an absolute constant. In the small rectangle
$0\le\Re(s-s_0)\le\delta$, $|\Im(s-s_0)|\le\delta$,
the other factors show $|\zeta(s)|\ll_C|G(s)|$.

The harmonic-measure kernel for the left boundary of the strip of width
one is
\[
 K_u(v)=\frac{\sin(\pi u)}{2(\cosh(\pi v)-\cos(\pi u))},
 \qquad 0<u<1.
\]
This formula follows by mapping $0<\Re z<1$ to the upper half-plane
by $e^{\pi i z}$ and changing variables in the Poisson kernel.
It is positive, has integral $1-u\le1$, and is bounded by
$C u/(u^2+v^2)$ for $u\le1/2$. Subharmonicity of $|G|^p$ and the
maximum principle on finite rectangles, followed by a limit, give
\[
 |\zeta(1/2+u+i\tau)|^p
 \ll_{p,C}1+\int_{\mathbb R}K_u(\tau-v)
 |\zeta(1/2+iv)|^p T^{-pC(v-t)^2}\,dv
\]
in the small rectangle. The horizontal boundary terms vanish because
of Gaussian decay and the polynomial bound just obtained.

Apply the disk submean inequality to $|\zeta|^p$ at $s_0$ with radius
$\delta$. Reflect the left half using \eqref{eq:localreflection},
and enlarge the right half to the small rectangle. For every $v$,
\[
 \int_0^\delta\int_{t-\delta}^{t+\delta}K_u(\tau-v)\,d\tau\,du
 \le\delta.
\]
After division by the disk area, the bound is
\[
 |\zeta(s_0)|^p\ll_p
 1+\delta^{-1}\int_{\mathbb R}
 |\zeta(1/2+iv)|^pT^{-pC(v-t)^2}\,dv.
\]
For $|v-t|>1/3$, the polynomial bound and sufficiently large $C_p$
make this integral $O(\delta)$, including $|v|$ far larger than $T$.
For the remaining interval discard the Gaussian. This proves
\eqref{eq:baselocal}. The disk submean inequality itself follows from
the circle mean inequality for the subharmonic function $|\zeta|^p$
by integrating the radius, including disks containing zeros.
\end{proof}

\begin{proposition}[fourth and twelfth moments from one sparse input]\label{prop:momentbaseproof}
Assuming only the sparse-fourth-moment estimate
\eqref{eq:basefourth} among the deep zeta moment inputs, one has
\[
 M_4(T)\prec T,\qquad M_{12}(T)\prec T^2,\qquad
 \cN(T,V)\prec\min(TV^{-4},T^2V^{-12})\quad(V\ge1).
\]
\end{proposition}
This derivation obtains the classical twelfth moment from the
sparse fourth moment. Attribution remains with the original authors;
this implication is Iwaniec's route~\cite{Iwa80}.
\begin{proof}
Take $J=1,L=T$ in \eqref{eq:basefourth} and sum dyadic height intervals
to obtain $M_4(T)\prec T$. Proposition~\ref{prop:localzetaproof}
with $p=4$ and disjoint neighborhoods of a $1$-separated set gives
$RV^4\prec T$; the added constant is absorbed for $V$ sufficiently
large, and the bounded-$V$ range is trivial.

Fix $\rho>0$, with all occurrences of $\prec$ in this proof taken at
a loss much smaller than $\rho$. If $V^4\le T^{1/2+\rho}$, the fourth
moment gives
$R\prec TV^{-4}\le T^{2+2\rho}V^{-12}$.
Otherwise put $L=cV^4T^{-\rho}$, with a fixed positive constant $c$.
If necessary increase the lower threshold by a constant so that
$L\ge\sqrt T$; the transition is covered by the preceding case.
When $R\ge1$, the just-proved fourth-moment estimate and the local
inequality give $V^4\ll T^{1+\rho/2}$, hence $L<T$ for large $T$.

Cover the disjoint neighborhoods $[t_j-1/3,t_j+1/3]$ by grid intervals
of length $L$. At most $2R$ grid intervals are needed. Boundary
intervals at $T,2T$ may be enlarged or shifted; a fixed number of
families in comparable dyadic height ranges handles them. Thus
\eqref{eq:basefourth} and the local lower mass yield
\[
 RV^4\prec RL+T\sqrt{R/L}.
\]
The $RL$ term is absorbed, since $L\ll V^4T^{-\rho}$.
Squaring the remaining inequality gives
$R\prec T^2/(LV^8)\ll T^{2+\rho}V^{-12}$.
As $\rho$ is arbitrary this proves the discrete twelfth-moment bound.

Finally the measure of each amplitude layer is at most a constant times
the maximal $1$-separated cardinality: choose a maximal separated
subset, whose unit neighborhoods cover that layer. Integrate in
dyadic $V$ and use the polynomial pointwise bound from the preceding
proof to ensure only $O(\log T)$ layers. This gives $M_{12}(T)\prec T^2$.
\end{proof}

\section{The derivative estimate from a precise mean-value input}\label{app:hbderivative}

We rederive the Heath-Brown estimate used here, keeping its attribution
to \cite{HB17}. Its sole nonelementary input in this appendix is
\begin{equation}\label{eq:vmvtinput}
 J_{s,q}(H):=\int_{[0,1]^q}\left|\sum_{h=1}^{\lfloor H\rfloor}
 e(\alpha_1h+\cdots+\alpha_qh^q)\right|^{2s}d\boldsymbol\alpha
 \ll_{q,\epsilon}H^{s+\epsilon},\qquad s=q(q+1)/2.
\end{equation}
This is the optimal Vinogradov mean-value theorem \cite{BDG16,Woo16}.

\begin{lemma}[coefficient clustering]\label{lem:hbcluster}
Suppose $d\ge3$ and $\lambda\le f^{(d)}\le A\lambda$ on an interval
of length $M$, with $A$ fixed. Put $H=\lfloor(A\lambda)^{-1/d}\rfloor$.
If $A\lambda$ is a sufficiently small constant depending on $d$, then
\[
 \mathcal P:=\#\left\{m,n:\left\|\frac{f^{(j)}(m)-f^{(j)}(n)}{j!}\right\|
 \le2H^{-j}\ (1\le j<d)\right\}
 \ll_{A,d}(M+\lambda M^2+\lambda^{-2/d})\log(2M).
\]
\end{lemma}
\begin{proof}
Keep the last two derivative conditions to obtain an upper bound. Let
$g_1=f^{(d-2)}/(d-2)!$, $g_2=f^{(d-1)}/(d-1)!$, and periodize
triangular functions of widths $B=4H^{2-d}$ and $C=4H^{1-d}$.
Their product $P(x,y)$ is at least $1/4$ on the counted differences,
vanishes unless $\|x\|\le B,\|y\|\le C$, and has nonnegative Fourier
coefficients $BC\,\operatorname{sinc}(rB)^2\operatorname{sinc}(sC)^2$.
Smallness of $\lambda$ ensures $B,C<1/2$.

Partition the integer interval into $K\asymp1+A\lambda M$ consecutive
pieces whose lengths make $|g_2(m)-g_2(n)|<1/2$ within each piece.
Fourier expansion and $|\sum_{i\le K}z_i|^2\le K\sum_i|z_i|^2$ show
\[
 \mathcal P\ll K\sum_i\sum_{m,n\text{ in piece }i}
 P(g_1(m)-g_1(n),g_2(m)-g_2(n)).
\]
Within a piece the $g_2$ condition is an ordinary real difference,
and monotonicity gives $|m-n|\ll\lambda^{-1}H^{1-d}\asymp H$.
The diagonal costs $M$. For $m=n+h$, $1\le h\ll H$, the derivative
of $g_1(x+h)-g_1(x)$ lies between fixed positive multiples of
$\lambda h$. A monotone function with derivative in $[\mu,A'\mu]$
takes integer arguments within distance $b$ of an integer at most
$O((1+\mu M)(1+b/\mu))$ times: count the integers in its range and
the length of each inverse interval. Hence the off-diagonal
count, even after dropping the piece restriction, is at most
\[
 \sum_{h\ll H}(1+\lambda Mh)(1+B/(\lambda h))
 \ll H+\lambda MH^2+(B/\lambda)\log(2H)+MBH
 \ll M+H^2\log(2M).
\]
Here $\lambda\asymp H^{-d}$ and $d\ge3$; if $H>M$ the trivial
$M^2\le H^2$ bound suffices. Multiply by $K$ and use
$\lambda M H^2\ll M$. This proves the bound.
\end{proof}

\begin{proposition}[Heath-Brown from \eqref{eq:vmvtinput}]
\label{prop:hbderivativeproof}
The mean-value bound \eqref{eq:vmvtinput} with $q=d-1$ implies
\eqref{eq:basederivative} for $d\ge3$.
\end{proposition}
\begin{proof}
Write $s=d(d-1)/2$ and first consider an integer interval of length $M$.
Conjugation makes $f^{(d)}$ positive. If $\lambda$ is bounded below,
the claimed first term bounds the sum trivially. If $H\ge M$, the
third claimed term does so. Otherwise $1\ll H<M$ and averaging
the $H$ translates of the sum, with end intervals costing $O(H)$, gives
\[
 \left|\sum_ne(f(n))\right|
 \ll H+H^{-1}\sum_n\max_{x\le H}
 \left|\sum_{h\le x}e\left(\sum_{j=1}^{d-1}\frac{f^{(j)}(n)}{j!}h^j\right)\right|.
\]
Taylor's remainder was removed by Abel summation: its derivative is
$O_d(A\lambda H^{d-1})=O_d(H^{-1})$.
To avoid a point-dependent maximum, use Fourier completion on
$\mathbb Z/(2H)\mathbb Z$, as in Lemma~\ref{lem:completionproof}.
This costs $O(\log(2H))$ and adds a single common linear coefficient
to every polynomial, preserving all coefficient differences.

Let $\nu(\boldsymbol\alpha)$ count the $n$ for which all the polynomial
coefficients are within $H^{-j}$ of $\alpha_j$ modulo one. Perturbing
coefficients by these amounts has total phase variation $O_d(1)$
on $[0,H]$. Abel summation and averaging over the coefficient box
therefore bound the preceding $n$ sum by
\[
 C_dH^s\left(I(H)+H^{-1}\int_0^H I(x)\,dx\right),\qquad
 I(x)=\int_{[0,1]^{d-1}}\nu(\boldsymbol\alpha)
 \left|\sum_{h\le x}e(\alpha_1h+\cdots+\alpha_{d-1}h^{d-1})\right|
 d\boldsymbol\alpha.
\]
Box volumes and overlaps give
$\int\nu\ll MH^{-s}$ and $\int\nu^2\ll\mathcal P H^{-s}$.
H\"older with three factors, of reciprocal exponents
$1-1/s$, $1/(2s)$ and $1/(2s)$, gives
\[
 I(x)\le(\int\nu)^{1-1/s}(\int\nu^2)^{1/(2s)}
 J_{s,d-1}(H)^{1/(2s)}.
\]
The monotonicity in $H$ follows directly by expanding the mean value
as a number of solutions. Applying \eqref{eq:vmvtinput} and collecting
the powers of $H$, which cancel, yields
\[
 \left|\sum_ne(f(n))\right|
 \ll_{d,A,\epsilon}H+M^{1-1/s+\epsilon}\mathcal P^{1/(2s)}.
\]
Insert Lemma~\ref{lem:hbcluster}. Since $2s=d(d-1)$, this is
\[
 \ll M^\epsilon\left(\lambda^{-1/d}
 +M^{1-1/[d(d-1)]}+M\lambda^{1/[d(d-1)]}
 +M^{1-2/[d(d-1)]}\lambda^{-2/[d^2(d-1)]}\right).
\]
If $\lambda<M^{-d}$ the last term is at least $M$. Otherwise
$\lambda^{-1/d}\le M$, and the last displayed term bounds it. It can therefore be removed in both cases. Substitute
$\lambda\asymp FM^{-d}$ to obtain \eqref{eq:basederivative}.
For a shorter subinterval first use its own length in this argument;
each term is increasing in that length at fixed $\lambda$, so replacing
it by $M$ is valid. The argument uses the phase on its given interval.
\end{proof}


\section{Theorems used}\label{app:base}

The preceding appendices derive the parts of \cite{B38} used here and
the reductions of de Faveri and Pandey. The analytic input theorems
and the arguments supplied here are identified individually below.

\subsection{Local zeta values and global moments}
Proposition~\ref{prop:localzetaproof} proves the local zeta estimate
\eqref{eq:baselocal}, including the functional-equation and growth
facts used in the proof. Proposition~\ref{prop:momentbaseproof} deduces
both $M_4(T)\prec T$ and $M_{12}(T)\prec T^2$ from the single sparse
fourth-moment input \eqref{eq:basefourth}. This cited theorem gives
the common analytic input for both global moments.

\subsection{Atkinson's short second moment}
With $T^\epsilon\le G\le T^{1/2-\epsilon}$ and
$S_t(\Delta,N)=\sum_{N<n\le N+\Delta}(-1)^nd(n)e(f(t,n))$, the input is
\begin{equation}\label{eq:baseatkinson}
 \begin{split}
 \int_{t-G}^{t+G}|\zeta(1/2+iu)|^2du
 &\ll G\log T\\
 &\quad+G\sum_{\substack{T^{1/3}\le N\ll T G^{-2}(\log T)^{1+\epsilon}\\
                         N\text{ dyadic}}}
 (TN)^{-1/4}\max_{0\le\Delta\le N}|S_t(\Delta,N)|.
 \end{split}
\end{equation}
This is the upper-bound consequence of \cite[Lemma 7.1]{dFP}.
Proposition~\ref{prop:atkinsondirect} proves it directly from the
double Dirichlet series, Euler summation, and Gaussian smoothing.
The estimate uses the exact phase at the center $t$ and includes all
Fourier indices and both nonstationary endpoints. For $t\asymp T$, changing
between the two comparable scales in this display costs only fixed
constants in the dyadic range.

\subsection{Sparse fourth moments}
For disjoint intervals $I_1,\ldots,I_J\subset[T,2T]$ of common length
$\sqrt T\le L\le T$, the input used in the pair count is
\begin{equation}\label{eq:basefourth}
 \sum_{j\le J}\int_{I_j}|\zeta(1/2+it)|^4dt
 \prec JL+T\sqrt{J/L}.
\end{equation}
It is Iwaniec's theorem \cite[Theorem~4]{Iwa80}, as stated in \cite[Proposition 8.1]{dFP}.
Proposition~\ref{prop:pairproof} derives the pair-count consequence
of this cited theorem.

\subsection{Heath-Brown's pair and Vinogradov mean values}
The pair in Lemma~\ref{lem:hb} is Theorem 2 of \cite{HB17}, following
there from Theorems 1 and 3. The underlying derivative estimate is
\begin{equation}\label{eq:basederivative}
 \left|\sum_{m\in I}e(g(m))\right|
 \ll_{d,\epsilon}M^{1+\epsilon}
 \left(F^{1/[d(d-1)]}M^{-1/(d-1)}
       +M^{-1/[d(d-1)]}+F^{-2/[d^2(d-1)]}\right),
\end{equation}
when $|g^{(d)}|\asymp FM^{-d}$ on an interval $I\subset(M,2M]$.
Proposition~\ref{prop:hbderivativeproof} proves this estimate from the
precise optimal Vinogradov mean-value input \eqref{eq:vmvtinput}.
The coefficient-box count, localization, all powers of the shift length
and the endpoint terms are included in that proof. The mean-value
input is the optimal Vinogradov mean-value theorem \cite{BDG16,Woo16}.

The finite-derivative conditions can then be settled explicitly
for the pair needed in this paper. For $F\ge M^{85/8}$ the target
$M^{1-a}F^{1/245}$ is at least $M$, so the trivial bound suffices.
For $M\le F\le M^2$, Weyl differencing and
Lemma~\ref{lem:secondproof}, with differencing length
$\asymp M F^{-1/3}$, give
$O(M^{1/2}F^{1/6}+M F^{-1/6})$; both terms are smaller than the
target. The differencing inequality follows by averaging the $H$
translates of the sum and applying Cauchy; in its usual form it is
\[
 |S|^2\ll M^2/H+(M/H)\sum_{h\le H}
 \left|\sum_m e(g(m+h)-g(m))\right|.
\]
For the reciprocal perturbations here the differenced second derivative
has size $hF/M^3$, so substitution proves the stated bound, including
end intervals. Constants in $H$ can be fixed to keep $m+h$ on a fixed
enlargement.

For $2\le\tau=\log F/\log M\le85/8$, apply
\eqref{eq:basederivative} with the integer $d\in\{3,\ldots,12\}$
whose interval
\[
 d-2+2/d\le\tau\le d-1+2/(d+1)
\]
contains $\tau$. The intervals are consecutive.
On it the maximum of the three relative exponents in
\eqref{eq:basederivative} is at most the chord joining
$-1/[d(d-1)]$ and $-1/[d(d+1)]$ at its two endpoints.
Its slope is $2/((d-1)^2(d+2))$. These slopes decrease with $d$,
so the resulting continuous piecewise linear function is concave and
lies below the line of its $d=8$ segment. That line is
\[
 \frac{\tau}{245}+\frac{269}{280}-\frac1{245}-1.
\]
This proves the desired exponent-pair estimate from
\eqref{eq:basederivative}, using derivatives only through order 12.
The exact phase estimates in Proposition~\ref{prop:exactphaseproof}
and Proposition~\ref{prop:tuplemorseproof} give them with relative
errors $O_d(|\eta|)$ through the specified finite derivative order.

The sixth-moment proof uses \eqref{eq:basefourth} and
\eqref{eq:vmvtinput}. Mean values of degree at most 11 cover the
chosen exponent pair.

The additional estimates in \cite{Supp} use direct differencing for
an alternative eighth-moment bound. The sixth-moment proof here uses
the two stated input theorems.
